\documentclass[10pt,a4paper]{amsart}
\usepackage[margin=19mm]{geometry}
\usepackage{amsmath,amssymb,mathtools}
\usepackage[scr=rsfs,cal=euler]{mathalfa}
\usepackage{xfrac}
\usepackage[unicode,hidelinks,bookmarks=false]{hyperref}
\newcommand{\ie}{i.e.\ }
\newcommand{\cf}{cf.\ }
\newcommand{\resp}{respectively\ }
\newcommand{\Subsection}{\subsection}
\providecommand{\nobreakdash}{\nobreak\hbox{-}}
\setlength{\emergencystretch}{2em}
\multlinegap0pt
\usepackage{amssymb}
\usepackage{csquotes}
\usepackage{algorithm}\usepackage{algpseudocode}
\newtheorem{theorem}{Theorem}
\title{An Extension of the Collatz Conjecture modulo $2^p+2^q$}
\author{Abderrahman Bouhamidi}
\date{Source: arXiv:2601.06208v1 (CC BY 4.0)}
\begin{document}
\maketitle

\noindent\emph{Source and licence.} An Extension of the Collatz Conjecture modulo $2^p+2^q$. Version arXiv:2601.06208v1 (CC BY 4.0), distributed by arXiv under the Creative Commons Attribution 4.0 International licence, http://creativecommons.org/licenses/by/4.0/, as recorded in the arXiv licence metadata for this exact version and shown on the abstract page https://arxiv.org/abs/2601.06208v1. Full licence notice: \texttt{licenses/arxiv-2601.06208-CC-BY-4.0.txt}.\par\smallskip

\noindent\textit{Editorial scope.} Counted unit: the article's theorem r)-(d,d+1,1,kappa0) (source0.tex lines 626--649). The article's complete original proof of it is delivered with it (source0.tex lines 650--683, one \texttt{proof} environment of 34 lines). No mathematics is written, repaired or re-derived: the statement and the proof are the article's own lines, in the article's own notation, and no retained line is substituted or re-flowed.  No further source-local context is needed: every cross-reference of every retained line resolves inside the two delivered spans.  Nothing was omitted from any retained span, so the package carries no omission record.  No retained line contains a citation, so the wrapper carries no bibliography and no bibliography entry.  No retained line contains a figure environment, an image-inclusion command or drawing code, so no image is delivered and the package declares no asset dependency.  Editorial formatting only, in addition: an AMS \texttt{amsart} preamble that restates the article's own declarations and macros used by the retained lines, namely the environments \texttt{theorem} on separate counters, together with the standard packages those lines need.  The retained environments print in this document as Theorem 1 in order of appearance, whereas the article prints the same items with the numbering of its own full document.  The complete native source of the article as distributed in the arXiv e-print for this exact version is bundled unchanged as \texttt{sources/192395/source0.tex}.
\begin{theorem} 
	Let $d\geq 2$ and consider     an admissible  triplet  $(d,\alpha,\beta)_{\pmb{\pm}}$  corresponding to   $\alpha=d+1$ and  $\beta=-\kappa_0$  with $\kappa_0=\pm 1$. 	Then, 	for all integers $a\geq 1$, for all integer $r\in\{1,\ldots, d-1\} $ and  for all integer $k\geq 0$, we have  the following relations 
	\begin{equation}\label{T(adk+kappa0 r)-(d,d+1,1,kappa0)}
		T^{(k)}(ad^k+\kappa_0 r)=a\alpha^k+\kappa_0r,
	\end{equation} 	
	and for $k>2$, we have
	\begin{equation}\label{T(adk+kappa0 (2d-1))-(d,d+1,1,kappa0)}
		T^{(k)}(ad^k+\kappa_0 (2d-1))=\left\{\begin{array}{lll}
			a3^{k-2} +\kappa_0 &\mathrm{if}& d=2,\\
			a\alpha^{k-1}+2\kappa_0 &\mathrm{if}& d>2,\\
		\end{array}\right.
	\end{equation}
	It follows immediately  that for $k\geq 0$:
	\begin{equation}\label{sigma(adk+kappa0 r)-(d,d+1,1,kappa0)}
		\sigma_\infty(ad^k+\kappa_0 r)=\sigma_\infty(a\alpha^k+\kappa_0 r)+k,
	\end{equation} 	
	and  for $k>2$:
	\begin{equation}\label{sigma(adk+kappa0 (2d-1))-(d,d+1,1,kappa0)}
		\sigma_\infty(ad^k+\kappa_0 (2d-1))=\left\{\begin{array}{lll}
			\sigma_\infty(a3^{k-2}+\kappa_0)+k &\mathrm{if}& d=2,\\
			\sigma_\infty(a\alpha^{k-1}+2\kappa_0)+k  &\mathrm{if}& d>2.\\
		\end{array}\right.
	\end{equation}
\end{theorem}

\begin{proof}
	\begin{itemize}
		\item For the relation  \eqref{T(adk+kappa0 r)-(d,d+1,1,kappa0)}, we again proceed by induction on $k\geq 0$. For $k=0$, the relation is trivially true. For $k=1$, we have 
		$T(ad+\kappa_0 r)=\dfrac{\alpha(ad+\kappa_0 r )+\beta r}{d}=a \alpha +\kappa_0 r  \dfrac{\alpha +\kappa_0\beta }{d}$.
		As  $\alpha +\kappa_0 \beta= d$. Then
		\begin{equation}\label{recurencek=1C} 
			T(ad+\kappa_0  r)=a \alpha +\kappa_0r, 
		\end{equation}	
		which shows that the formula is also true for $k=1$. We assume that, for all $a\geq 1$  and for all $r\in\{1,\ldots,d-1\}$,
		the formula \eqref{T(adk+kappa0 r)-(d,d+1,1,kappa0)} is true until an order $k$.   So,  by  the recurrence hypothesis at step $k$, we get
		$$T^{(k+1)}(ad^{k+1}+\kappa_0 r)=T\Bigl[T^{(k)}((ad)d^{k}+\kappa_0 r )\Bigr]=T\Bigl[(ad)\alpha^{k}+\kappa_0r\Bigr].$$
		According to the relation \eqref{recurencek=1C} at step 1, we obtain
		$T^{(k+1)}(ad^{k+1}+\kappa_0 r)=a\alpha^{k+1}+\kappa_0r$.
		Which show that the formula  \eqref{T(adk+kappa0 r)-(d,d+1,1,kappa0)}  is true at the order $k+1$. The relation \eqref{sigma(adk+kappa0 r)-(d,d+1,1,kappa0)} is a consequence of the relation  \eqref{T(adk+kappa0 r)-(d,d+1,1,kappa0)} .
		
		\item For the relation  \eqref{T(adk+kappa0 (2d-1))-(d,d+1,1,kappa0)}, for $k> 2$ we have
		$$T^{(k)}(ad^k+\kappa_0 (2d-1))=T^{(k-1)}\bigl[T(ad^k+\kappa_0 (2d-1))\bigr].$$
		As  $\bigl[\kappa_0(ad^k+\kappa_0 (2d-1))\bigr]_d=d-1$ and $\alpha-1=d$, it follows 
		$$T(ad^k+\kappa_0 (2d-1))=\dfrac{\alpha(ad^k+\kappa_0 (2d-1))-\kappa_0(d-1)}{d}=\alpha ad^{k-1}+2\kappa_0d.$$
		Then,  we have
		\begin{equation}\label{R-3}
			T^{(k)}(ad^k+\kappa_0 (2d-1))=\left\{\begin{array}{lll}
				T^{(k-3)}\bigl[(3 a)2^{k-3}+\kappa_0\bigr] &\mathrm{if}& d=2,\\
				T^{(k-2)}\bigl[(\alpha a)d^{k-2}+2\kappa_0\bigr] &\mathrm{if}& d>2,\\
			\end{array}\right.
		\end{equation} 
		Then, using the relation    \eqref{T(adk+kappa0 r)-(d,d+1,1,kappa0)}, we obtain the relation \eqref{T(adk+kappa0 (2d-1))-(d,d+1,1,kappa0)}.
		The relations \eqref{sigma(adk+kappa0 r)-(d,d+1,1,kappa0)}
		and  \eqref{sigma(adk+kappa0 (2d-1))-(d,d+1,1,kappa0)}
		are  immediate consequences of the relations 
		\eqref{T(adk+kappa0 r)-(d,d+1,1,kappa0)} and
		\eqref{T(adk+kappa0 (2d-1))-(d,d+1,1,kappa0)}, respectively.
	\end{itemize}
	\hfill\end{proof}

\end{document}
